\section*{Erd\H{o}s Problem 801: a small set with many edges}

\paragraph{Statement and claim.}
For every sufficiently large \(n\), if an \(n\)-vertex graph
has independence number at most \(\sqrt n\), then some
set of at most \(\sqrt n\) vertices spans
\(\Omega(\sqrt n\log n)\) edges.
We reconstruct Alon's proof, allowing the non-strict
independence hypothesis in the displayed question:
\url{https://web.math.princeton.edu/~nalon/PDFS/indrams1.pdf},
Section 3. See \url{https://www.erdosproblems.com/801}.
The external input is the Ajtai--Koml\'os--Szemer\'edi
triangle-free independence estimate: a triangle-free
\(v\)-vertex graph of average degree \(d\geq2\) has an
independent set of size at least \(c(v/d)\log d\),
for an absolute \(c>0\).

\paragraph{Two direct sampling cases.}
Put \(m=\lfloor\sqrt n\rfloor\), and let \(\bar d\) be the
average degree of \(G\). If \(\bar d\geq\sqrt n\log n\),
a uniform \(m\)-vertex subset spans an expected
\[
 e(G)\frac{\binom m2}{\binom n2}
 =\Omega(\sqrt n\log n)
\]
edges, so one such subset suffices.
Otherwise at least \(n/2\) vertices have degree at most
\(2\sqrt n\log n\). Let \(G'\) be their induced graph,
so its maximum degree obeys that bound.

Suppose a neighbourhood in \(G'\) spans at least
\(\sqrt n(\log n)^3\) edges.
If the neighbourhood has at most \(m\) vertices it already
works. If it is larger, its order is at most
\(2\sqrt n\log n\), and a uniform \(m\)-subset has expected
edge count at least a constant times
\[
 \sqrt n(\log n)^3
 \frac{n}{n(\log n)^2}
 =\sqrt n\log n.
\]
Again the assertion follows.

\paragraph{Removing triangles after random sampling.}
We may now assume every neighbourhood in \(G'\) spans
fewer than \(\sqrt n(\log n)^3\) edges.
Each triangle is counted in three neighbourhoods, so
\(G'\) has at most \(n^{3/2}(\log n)^3\) triangles.
Independently retain each vertex with probability
\(p=n^{-2/5}\).
The expected retained vertex count minus the retained
triangle count is at least
\[
 \frac n2p-n^{3/2}(\log n)^3p^3
 =\frac12n^{3/5}-n^{3/10}(\log n)^3
 \geq\frac14 n^{3/5}
\]
for large \(n\).
Some sample attains this lower bound.
Remove one vertex from each of its triangles, and then
discard extra vertices if needed.
We obtain a triangle-free induced subgraph \(J\) of order
\(v=\lfloor n^{3/5}/4\rfloor\).

\paragraph{The independence bound forces enough edges.}
Write \(d\) for the average degree of \(J\).
First, the elementary bound \(\alpha(J)\geq v/(d+1)\)
and \(\alpha(J)\leq\alpha(G)\leq\sqrt n\) show
\[
 d+1\geq v/\sqrt n=\Omega(n^{1/10}).
\]
Thus \(d\geq2\) and \(\log d\geq c_1\log n\) for large \(n\).
The cited triangle-free theorem then gives
\[
 \sqrt n\geq\alpha(J)\geq c\frac vd\log d,
 \qquad d\geq c_2n^{1/10}\log n.
\]
Finally choose a uniform \(m\)-subset of \(V(J)\).
Its expected edge count is
\[
 \frac{vd}{2}\frac{\binom m2}{\binom v2}
 =\Omega\left(\frac{dm^2}{v}\right)
 =\Omega(\sqrt n\log n).
\]
Some subset therefore has the required number of edges.

\paragraph{Scope.}
All reductions, sampling arguments, and the transfer back
to a set of at most \(\sqrt n\) vertices are proved above.
The full assertion uses only the explicitly stated
triangle-free independence theorem as a substantial
external input; the conclusion is not claimed for tiny
\(n\) where the asymptotic notation has no intended force.

\paragraph{Completion estimate.}
\textbf{COMPLETION: 100\%}
